\documentclass[a4paper,10pt]{article}
\usepackage[T1]{fontenc}
\usepackage{lmodern,microtype}
\usepackage[margin=25mm,headheight=14pt]{geometry}
\usepackage{amsmath,amssymb,amsthm,mathtools}
\usepackage{xcolor,enumitem,booktabs,fancyhdr}
\definecolor{navy}{HTML}{123B63}
\usepackage[colorlinks=true,linkcolor=navy,urlcolor=navy,citecolor=navy]{hyperref}
\hypersetup{pdftitle={Bootstrap Arguments: From an ODE to a Simple PDE},pdfauthor={Jiten Kumbhakar},pdfsubject={Self-contained note on continuity bootstraps, local existence and continuation}}
\numberwithin{equation}{section}
\newtheorem{theorem}{Theorem}[section]
\newtheorem{lemma}[theorem]{Lemma}
\newtheorem{proposition}[theorem]{Proposition}
\theoremstyle{definition}
\newtheorem{definition}[theorem]{Definition}
\theoremstyle{remark}
\newtheorem*{remark}{Remark}
\newcommand{\R}{\mathbb R}
\newcommand{\T}{\mathbb T}
\newcommand{\norm}[1]{\lVert#1\rVert}
\newcommand{\abs}[1]{\lvert#1\rvert}
\newcommand{\dd}{\,\mathrm d}
\newcommand{\Tmax}{T_{\max}}
\setlength{\parindent}{0pt}
\setlength{\parskip}{5pt plus 1pt}
\setlist[enumerate]{leftmargin=*,itemsep=2pt,topsep=3pt}
\setlist[itemize]{leftmargin=*,itemsep=2pt,topsep=3pt}
\pagestyle{fancy}
\fancyhf{}
\fancyhead[L]{\small\color{navy}Bootstrap arguments}
\fancyhead[R]{\small Jiten Kumbhakar}
\fancyfoot[C]{\small\thepage}
\renewcommand{\headrulewidth}{0.3pt}
\newcommand{\pageheading}[1]{\section{#1}}
\begin{document}
\thispagestyle{plain}
\begin{center}
{\LARGE\bfseries\color{navy}Bootstrap Arguments}\\[5pt]
{\large From an ODE to a Simple PDE}\\[8pt]
Jiten Kumbhakar\\[3pt]
{\small Expository note \;\textbullet\; Version 1.0 \;\textbullet\; 23 October 2026}
\end{center}
\vspace{3pt}
\textbf{Purpose.} A reusable, self-contained introduction to the continuity bootstrap:
first the logical principle, then an ordinary differential equation, and finally a
semilinear heat equation. The local construction and the continuation step are
included. Only elementary calculus, uniform convergence and basic normed-space
notation are assumed; the fixed-point fact used below is proved in Section~2.

\pageheading{What is a bootstrap argument?}
\begin{definition}[Continuity bootstrap]
A continuity bootstrap proves a bound by assuming a weaker bound on an initial
time interval, deriving a \emph{strictly better} bound there, and using continuity
to show that the weaker bound can never first fail.
\end{definition}
The temporary assumption is called the \emph{bootstrap hypothesis}; the stronger
estimate is the \emph{bootstrap improvement}. For example,
\[
X\leq 2\varepsilon
\quad\Longrightarrow\quad
X\leq \tfrac32\varepsilon<2\varepsilon.
\]
This is not a circular proof: the initial data start inside the allowed region,
and the improvement rules out the first time the solution could reach its boundary.

\begin{lemma}[First-exit principle]\label{lem:exit}
Let $X:[0,\Tmax)\to[0,\infty)$ be continuous, where $0<\Tmax\leq\infty$.
Let $0\leq m<M$, and assume $X(0)<M$. Suppose that for every $t<\Tmax$,
\[
\bigl[X(s)\leq M\text{ for every }s\in[0,t]\bigr]
\quad\Longrightarrow\quad X(t)\leq m.
\]
Then $X(t)\leq m$ for every $t<\Tmax$.
\end{lemma}
\begin{proof}
If $X(t_1)\geq M$ for some $t_1<\Tmax$, continuity gives a first time
$\tau\in(0,t_1]$ with $X(\tau)=M$. On $[0,\tau]$ we have $X\leq M$, so the
assumed improvement gives $X(\tau)\leq m<M$, a contradiction. Thus $X<M$
everywhere. Applying the improvement on each $[0,t]$ now gives $X(t)\leq m$.
\end{proof}

\textbf{The usual quantity.} If $f:[0,\Tmax)\to[0,\infty)$ is continuous, set
$X(t)=\max_{0\leq s\leq t}f(s)$. This running maximum is continuous and
nondecreasing. Indeed, on any compact time interval let $\omega$ be a modulus
of continuity of $f$. For $t\geq r$, the new values satisfy
$f(s)\leq f(r)+\omega(t-r)\leq X(r)+\omega(t-r)$ for $r\leq s\leq t$.
Hence $0\leq X(t)-X(r)\leq\omega(t-r)$. A hypothesis $X(t)\leq M$ therefore
already controls all earlier times.

\textbf{Scope.} Here ``bootstrap'' means a bound closed by continuity. The related
\emph{regularity bootstrap} repeatedly improves differentiability using the
equation; it is a different use of the same word.

\newpage
\pageheading{The estimate that closes the argument}
\begin{proposition}[Quadratic bootstrap]\label{prop:quadratic}
Let $a>0$, $b\geq0$, and let $X$ be continuous and nonnegative on
$[0,\Tmax)$ with $X(0)\leq a$. Suppose that, whenever $X(s)\leq2a$ on
$[0,t]$, the estimate
\begin{equation}\label{eq:quadratic}
X(t)\leq a+bX(t)^2
\end{equation}
holds. If $8ab\leq1$, then $X(t)\leq\frac32a$ for every $t<\Tmax$.
\end{proposition}
\begin{proof}
Under the bootstrap hypothesis,
\[
X(t)\leq a+b(2a)^2=a+4ba^2\leq\tfrac32a<2a.
\]
Apply Lemma~\ref{lem:exit} with $M=2a$ and $m=\frac32a$.
\end{proof}
The numbers $2$ and $3/2$ provide convenient slack; they have no intrinsic
significance. More generally, $X\leq a+bX^p$, $p>1$, closes in exactly the
same way if $b(2a)^p\leq a/2$.

\textbf{Why the inequality alone is insufficient.}
For $b>0$, every $X\geq1/b$ satisfies $X\leq a+bX^2$. Thus
\eqref{eq:quadratic} allows a large branch as well as a small one. The initial
bound and continuity keep the solution on the small branch. When $a=0$,
the thresholds above coincide; in our examples the zero solution is instead
handled directly by uniqueness.

\textbf{Finite time versus global time.}
If $b=b(T)$ grows with the time horizon, the smallness condition only gives
control on horizons for which $8ab(T)\leq1$. For example, the integral form
of $y'=y^2$ gives $X(T)\leq\abs{y_0}+T X(T)^2$, with
$X(T)=\max_{s\leq T}\abs{y(s)}$. It yields a useful short-time estimate,
not a bound uniform for all time. The damping in our two examples removes
this growth of the nonlinear constant.

\subsection*{The elementary existence tool}
A normed space is \emph{complete} if every Cauchy sequence converges in it;
it is then called a Banach space. The space of continuous functions on a
compact interval, with the uniform norm, is complete: uniform limits of
continuous functions are continuous, and pointwise completeness gives the limit.
The same argument applies to continuous functions with values in a Banach space.

\begin{lemma}[Contraction principle]\label{lem:contraction}
Let $B$ be a nonempty closed subset of a Banach space. If
$\Phi:B\to B$ satisfies
$\norm{\Phi(v)-\Phi(w)}\leq q\norm{v-w}$ for some $0\leq q<1$, then
$\Phi$ has exactly one fixed point in $B$.
\end{lemma}
\begin{proof}
Starting with $v_0\in B$, let $v_{n+1}=\Phi(v_n)$. Induction gives
$\norm{v_{n+1}-v_n}\leq q^n\norm{v_1-v_0}$. Summing this geometric bound
shows that $(v_n)$ is Cauchy. Its limit $v\in B$ satisfies $v=\Phi(v)$ by
the Lipschitz estimate. If $w$ is another fixed point, then
$\norm{v-w}\leq q\norm{v-w}$, which forces $v=w$.
\end{proof}

\textbf{Division of work.} A fixed-point argument constructs a local solution.
A bootstrap controls a solution on its existing interval. A continuation
criterion determines whether that interval can be extended. Each step has
its own hypothesis; a bootstrap bound alone is not an existence theorem.

\newpage
\pageheading{ODE example: damping against quadratic growth}
Consider the real-valued initial-value problem
\begin{equation}\label{eq:ode}
y'(t)+y(t)=y(t)^2,\qquad y(0)=y_0\in\R.
\end{equation}
\begin{theorem}\label{thm:ode}
If $\delta=\abs{y_0}\leq1/8$, problem~\eqref{eq:ode} has a unique global
$C^1$ solution and
\begin{equation}\label{eq:odebound}
\abs{y(t)}\leq \tfrac32\delta e^{-t}\leq2\delta e^{-t}
\qquad(t\geq0).
\end{equation}
\end{theorem}
\begin{proof}
\textbf{1. Local construction and continuation.}
Multiplication by $e^t$ gives the equivalent integral equation
\begin{equation}\label{eq:odeintegral}
y(t)=e^{-t}y_0+\int_0^t e^{-(t-s)}y(s)^2\dd s.
\end{equation}
Let $A=\max\{1,\abs{y_0}\}$ and $h=1/(8A)$. On the closed ball
$\norm{v}_{C([0,h])}\leq2A$, the right-hand side defines $\Phi$ with
\[
\norm{\Phi(v)}\leq A+4A^2h\leq\tfrac32A,\qquad
\norm{\Phi(v)-\Phi(w)}\leq4Ah\norm{v-w}\leq\tfrac12\norm{v-w}.
\]
Lemma~\ref{lem:contraction} gives a fixed point, and differentiation of
\eqref{eq:odeintegral} shows that it solves~\eqref{eq:ode}. Two bounded
solutions coincide on sufficiently short intervals by the same difference
estimate; subdivision gives uniqueness on their common interval.

Restarting at any time produces a maximal solution on $[0,\Tmax)$.
If $\Tmax<\infty$ and $\sup_{t<\Tmax}\abs{y(t)}\leq B<\infty$, the same
construction from any initial time has lifespan at least
$1/(8\max\{1,B\})$. Restarting sufficiently close to $\Tmax$ extends past
$\Tmax$, contradicting maximality. Thus a uniform bound prevents finite-time
breakdown.

\textbf{2. Bootstrap estimate.}
If $\delta=0$, uniqueness gives $y=0$. Otherwise define
\[
X(t)=\max_{0\leq s\leq t} e^s\abs{y(s)},\qquad t<\Tmax.
\]
This is continuous and $X(0)=\delta$. From~\eqref{eq:odeintegral}, for
each $r\leq t$,
\[
e^r\abs{y(r)}
\leq\delta+\int_0^r e^s\abs{y(s)}^2\dd s
\leq\delta+X(t)^2\int_0^r e^{-s}\dd s
\leq\delta+X(t)^2.
\]
Taking the maximum in $r$ gives $X(t)\leq\delta+X(t)^2$.
Since $8\delta\leq1$, Proposition~\ref{prop:quadratic} yields
$X(t)\leq3\delta/2$ for all $t<\Tmax$.

\textbf{3. Global continuation.}
Consequently $\abs{y(t)}\leq3\delta/2$ on the maximal interval. Step 1
forces $\Tmax=\infty$, and the weighted estimate gives~\eqref{eq:odebound}.
\end{proof}

\textbf{A check using the exact solution.}
Direct differentiation verifies
\[
y(t)=\frac{y_0e^{-t}}{1-y_0(1-e^{-t})}
\]
as long as the denominator is nonzero. This confirms that $1/8$ is only a
convenient sufficient threshold. If $y_0>1$, the solution blows up at
$t=\log(y_0/(y_0-1))$; if $y_0=1$, it stays equal to $1$. Smallness is
therefore meaningful, although our chosen constant is not sharp.

\newpage
\pageheading{PDE example: the setting and the heat operator}
Let $\T=\R/(2\pi\mathbb Z)$ be the circle, identified with $[0,2\pi]$
with the endpoints joined. Consider
\begin{equation}\label{eq:pde}
\partial_tu-\partial_x^2u+u=u^2,\qquad u(0,x)=u_0(x),\qquad x\in\T.
\end{equation}
All functions are real-valued and periodic in $x$. We work in
$C(\T)$, the Banach space of continuous periodic functions, with
$\norm{f}_\infty=\max_{x\in\T}\abs{f(x)}$. In particular,
\begin{equation}\label{eq:product}
\norm{f^2}_\infty\leq\norm{f}_\infty^2,\qquad
\norm{f^2-g^2}_\infty\leq
(\norm{f}_\infty+\norm{g}_\infty)\norm{f-g}_\infty.
\end{equation}
No Sobolev embedding or product theorem is needed in this example.

\subsection*{An explicit linear solution operator}
For $t>0$, define
\[
G_t(z)=\frac{1}{\sqrt{4\pi t}}\exp\!\left(-\frac{z^2}{4t}\right),\qquad
K_t(x)=\sum_{k\in\mathbb Z}G_t(x+2\pi k),
\]
and
\[
S(t)f(x)=\int_0^{2\pi}K_t(x-z)f(z)\dd z,\qquad
H(t)f=e^{-t}S(t)f,\qquad H(0)f=f.
\]
The Gaussian series and its derivatives converge uniformly when $t$ stays
away from zero. Since $G_t\geq0$ and $\int_\R G_t=1$, periodization gives
$K_t\geq0$ and $\int_0^{2\pi}K_t=1$. Consequently
\begin{equation}\label{eq:heatdecay}
\norm{S(t)f}_\infty\leq\norm{f}_\infty,
\qquad \norm{H(t)f}_\infty\leq e^{-t}\norm{f}_\infty.
\end{equation}
Completing the square in the Gaussian convolution gives
$G_t*G_s=G_{t+s}$; periodizing this identity gives
$H(t)H(s)=H(t+s)$. Such a family is called a \emph{semigroup}.

Also $H(t)f\to f$ uniformly as $t\downarrow0$. To see this, extend $f$
periodically to $\R$ and write
\[
S(t)f(x)-f(x)=\int_\R G_t(z)\bigl(f(x-z)-f(x)\bigr)\dd z.
\]
On $\abs z<\eta$, uniform continuity makes the difference small; on
$\abs z\geq\eta$, bound it by $2\norm f_\infty$ and use that the Gaussian
tail tends to zero. The factor $e^{-t}$ tends to one. The semigroup identity
and~\eqref{eq:heatdecay} then give continuity at every time.

\begin{definition}[Mild solution]\label{def:mild}
For $u_0\in C(\T)$, a mild solution of~\eqref{eq:pde} on $[0,T]$ is
$u\in C([0,T];C(\T))$ satisfying
\begin{equation}\label{eq:mild}
u(t)=H(t)u_0+\int_0^t H(t-s)u(s)^2\dd s.
\end{equation}
The integral is the norm limit of Riemann sums of a continuous
$C(\T)$-valued function. This formula incorporates the initial condition.
\end{definition}
For a smooth solution, differentiating $H(t-s)u(s)$ in $s$ gives
$H(t-s)u(s)^2$, so integration yields~\eqref{eq:mild}. Conversely the mild
formula solves~\eqref{eq:pde} in the weak sense: for every smooth periodic
$\varphi$, setting $a(t)=\int_\T u(t)\varphi\dd x$ gives
\[
a'(t)=\int_\T u(t)(\varphi''-\varphi)\dd x
       +\int_\T u(t)^2\varphi\dd x.
\]
Indeed, move the symmetric kernel $H$ onto $\varphi$ in~\eqref{eq:mild}
and differentiate; $\partial_tH(t)\varphi=H(t)(\varphi''-\varphi)$.
Thus the integral formulation has a precise PDE meaning without assuming
derivatives of $u$ in advance.

\newpage
\pageheading{Local PDE solutions and the continuation criterion}
\begin{proposition}\label{prop:local}
For every $u_0\in C(\T)$, there is a unique maximal mild solution of
\eqref{eq:pde} on $[0,\Tmax)$. If $\Tmax<\infty$, then
\begin{equation}\label{eq:blowup}
\sup_{0\leq t<\Tmax}\norm{u(t)}_\infty=\infty.
\end{equation}
\end{proposition}
\begin{proof}
\textbf{1. A fixed-point space.}
Put $A=\max\{1,\norm{u_0}_\infty\}$, $h=1/(8A)$, and
\[
E=C([0,h];C(\T)),\qquad
\norm v_E=\max_{0\leq t\leq h}\norm{v(t)}_\infty,\qquad
B=\{v\in E:\norm v_E\leq2A\}.
\]
Both $C(\T)$ and $E$ are complete, and $B$ is a nonempty closed subset.
Let $\Phi(v)$ be the right-hand side of~\eqref{eq:mild}, with $u$ replaced
by $v$. The continuity of $H$ proved above and its uniform bound ensure
$\Phi(v)\in E$. More explicitly, $v(s)^2$ is continuous, $H(t)$ is
strongly continuous, and their composition is jointly continuous for
$0\leq s\leq t\leq h$; integration therefore preserves continuity.

\textbf{2. Invariance and contraction.}
Using~\eqref{eq:product}--\eqref{eq:heatdecay}, for $v,w\in B$,
\begin{align*}
\norm{\Phi(v)}_E
&\leq\norm{u_0}_\infty+h\norm v_E^2
 \leq A+4A^2h=\tfrac32A<2A,\\
\norm{\Phi(v)-\Phi(w)}_E
&\leq h(\norm v_E+\norm w_E)\norm{v-w}_E
 \leq4Ah\norm{v-w}_E=\tfrac12\norm{v-w}_E.
\end{align*}
Lemma~\ref{lem:contraction} gives a mild solution on $[0,h]$.

\textbf{3. Uniqueness and restarting.}
If $u,v$ are two mild solutions with the same initial value, bounded by
$B_0$ on a compact common interval, then on a subinterval of length $\ell$
starting with equal values their difference satisfies
\[
\max\norm{u-v}_\infty\leq2B_0\ell\max\norm{u-v}_\infty.
\]
Choose $2B_0\ell<1$ (the case $B_0=0$ is immediate), then cover the
common interval by finitely many such subintervals. This proves uniqueness
among all continuous mild solutions, not only inside the construction ball.

The restart identity used here follows by splitting~\eqref{eq:mild} at
$t_0$ and using the semigroup property:
\begin{equation}\label{eq:restart}
u(t)=H(t-t_0)u(t_0)+\int_{t_0}^t H(t-s)u(s)^2\dd s,
\qquad t\geq t_0.
\end{equation}
Local solutions can therefore be joined consistently, giving a maximal
interval $[0,\Tmax)$.

\textbf{4. Continuation uses exactly the controlled norm.}
Suppose $\Tmax<\infty$ but $\norm{u(t)}_\infty\leq B_1$ for all
$t<\Tmax$. At every such $t_0$, Steps 1--2 give a solution starting from
$u(t_0)$ for at least the common duration
\[
h_* =\frac{1}{8\max\{1,B_1\}}>0.
\]
Take $t_0<\Tmax$ with $\Tmax-t_0<h_*$. Uniqueness identifies this
solution with the old one on the overlap, and it continues past $\Tmax$.
This contradicts maximality and proves~\eqref{eq:blowup}.
\end{proof}

Notice that no limit $u(\Tmax)$ had to be constructed. A local lifespan
depending only on the bounded norm allows us to restart \emph{before} the
putative final time. This is often the cleanest continuation argument.

\newpage
\pageheading{Closing the PDE bootstrap and reusing the method}
\begin{theorem}\label{thm:pde}
Let $u_0\in C(\T)$ and $\delta=\norm{u_0}_\infty\leq1/8$. Then
\eqref{eq:pde} has a unique global mild solution and
\begin{equation}\label{eq:pdebound}
\norm{u(t)}_\infty\leq\tfrac32\delta e^{-t}\leq2\delta e^{-t}
\qquad(t\geq0).
\end{equation}
No sign assumption is imposed on $u_0$.
\end{theorem}
\begin{proof}
Use Proposition~\ref{prop:local} to obtain the maximal solution. If
$\delta=0$, uniqueness gives $u=0$. For $\delta>0$, set
\[
X(t)=\max_{0\leq s\leq t}e^s\norm{u(s)}_\infty,\qquad t<\Tmax.
\]
The norm is continuous since
$\abs{\norm{u(t)}_\infty-\norm{u(r)}_\infty}\leq\norm{u(t)-u(r)}_\infty$;
the running-maximum argument of Section~1 makes $X$ continuous. Also
$X(0)=\delta$. For $r\leq t$, equations~\eqref{eq:mild} and
\eqref{eq:heatdecay} imply
\begin{align*}
e^r\norm{u(r)}_\infty
&\leq\delta+\int_0^r e^s\norm{u(s)^2}_\infty\dd s\\
&\leq\delta+X(t)^2\int_0^r e^{-s}\dd s
 \leq\delta+X(t)^2.
\end{align*}
Maximizing in $r$ gives $X(t)\leq\delta+X(t)^2$. Under $X\leq2\delta$,
the improvement is $X\leq\delta+4\delta^2\leq3\delta/2<2\delta$.
Proposition~\ref{prop:quadratic} therefore gives $X(t)\leq3\delta/2$ on
the entire maximal interval. In particular $\norm{u(t)}_\infty$ stays
bounded. The continuation criterion forces $\Tmax=\infty$, and the
weighted bound proves~\eqref{eq:pdebound}.
\end{proof}

\textbf{What changed from the ODE?} The scalar propagator $e^{-t}$ became
the operator $H(t)$, absolute value became the uniform norm, and the scalar
product estimate became~\eqref{eq:product}. The continuity argument itself
was unchanged. Spatially constant initial data give the ODE exactly, since
$S(t)$ preserves constants. Thus large constant data also show why this PDE
cannot have a small-data conclusion for arbitrary data.

\subsection*{A checklist for the next equation}
\begin{enumerate}
\item \textbf{Construct locally.} Specify the solution class and its maximal interval.
\item \textbf{Choose the quantity.} Define $X(t)$, possibly with a time weight;
prove its continuity and compute $X(0)$.
\item \textbf{State the temporary bound.} Choose a threshold strictly above
the initial size, e.g.\ $X\leq2a$.
\item \textbf{Improve with explicit constants.} Derive the estimate under
that hypothesis and arrange $X\leq3a/2$ by a stated smallness condition.
\item \textbf{Close by first exit.} Conclude the estimate for every
$t<\Tmax$, not yet for times at or beyond $\Tmax$.
\item \textbf{Continue in the correct norm.} Match the bound to the actual
continuation criterion. Check that the constants are uniform in the time horizon.
\end{enumerate}
An estimate that returns only $X\leq2a$ has no strict margin. A bound in a
weaker norm does not exclude breakdown in a stronger norm unless the local
theory supplies the needed continuation criterion.

\subsection*{Further reading}
All arguments needed for the stated results are given above. For wider context,
see T.~Tao, \emph{Nonlinear Dispersive Equations: Local and Global Analysis},
CBMS~106, AMS, 2006, \S1.3, ``Bootstrap and continuity arguments''
(\href{https://www.math.ucla.edu/~tao/preprints/chapter.pdf}{author's online text});
and J.~K.~Hunter, \emph{Notes on Partial Differential Equations}, Chapter~5,
``The heat equation''
(\href{https://www.math.ucdavis.edu/~hunter/pdes/ch5.pdf}{online lecture notes}).
\end{document}
